\chapter{一维定态问题}

\newpage


\textbf{1-1}\quad 一个质量为 $\displaystyle \mu$ 的粒子在一维势场 $\displaystyle V(x) = \begin{cases} 0, & |x| > a \\ -V_0, & |x| < a \end{cases}$ 中运动, 其中 $\displaystyle V_0 > 0$ . 求基态能量 $\displaystyle E_0$ 满足的方程; 求存在且仅存在一个束缚态的条件.

\vfill

\newpage

\textbf{1-2}\quad 质量为 $\displaystyle \mu$ 的粒子在势场 $\displaystyle V(x) = -\alpha \delta (x)(\alpha >0)$ 中运动, 求束缚态能级和相应的波函数.

\vfill

\newpage

\textbf{1-3}\quad 质量为 $\displaystyle \mu$ 的粒子处于一维势场

$$
V (x) = \left\{ \begin{array}{l l} \infty , & x \geqslant a \\ 0, & 0 <   x <   a \\ \infty , & x = 0 \\ 0, & - a <   x <   0 \\ \infty , & x \leqslant - a \end{array} \right.
$$

中, 求定态能量 E 与波函数 $\displaystyle \psi(x)$ .

\vfill

\newpage

\textbf{1-4}\quad 求在半壁无限深方势阱 $\displaystyle V(x) = \begin{cases} \infty, & x < 0 \\ 0, & 0 < x < a \\ V_0, & x > a \end{cases}$ 中存在束缚态的条件 ( $\displaystyle V_0 > 0$ ).

\vfill

\newpage

\textbf{1-5}\quad 质量为 $\displaystyle \mu$ 的粒子在一维势场

$$
V (x) = - \alpha \delta (x) + V ^ {\prime} (x), V ^ {\prime} (x) = \left\{ \begin{array}{l l} 0, & x <   0 \\ V _ {0}, & x > 0 \end{array} \right.
$$

中运动, 其中 $\displaystyle \alpha$ 与 $\displaystyle V_{0}$ 均为正实数. (1)给出存在束缚态的条件, 并给出能量本征值与相应的本征函数; (2)给出粒子处于 $\displaystyle x > 0$ 区中的概率, 它是大于 $\displaystyle 1/2$ , 还是小于 $\displaystyle 1/2$ ? 为什么?

\vfill

\newpage

\textbf{1-6}\quad 一个质量为 $\displaystyle \mu$ 的粒子在一维势场

$$
V (x) = \left\{ \begin{array}{l l} \infty , & x <   0, x > a \\ \alpha \delta (x - (a / 2)), & 0 <   x <   a \end{array} \right.
$$

中运动, 其中 $\displaystyle \alpha$ 与 $\displaystyle a$ 是正的常数. 求第一激发态能量, 并讨论 $\displaystyle \alpha \to 0$ 时的定态能量.

\vfill

\newpage

\textbf{1-7}\quad 在谐振子势的中心附加 $\displaystyle \delta$ 函数势, $\displaystyle V(x) = \dfrac{1}{2} \mu \omega^2 x^2 + V_0 \delta(x)$ . 原来谐振子定态解中, 哪些仍是解? 哪些不再是解? 要重新求.

\vfill

\newpage

\textbf{1-8}\quad 求质量为 $\displaystyle \mu$ 的粒子在势场 $\displaystyle V(x) = \begin{cases} -\dfrac{a}{x}, & x > 0 \\ \infty, & x < 0 \end{cases}$ 中的束缚定态能量与波函数, 其中 $\displaystyle a > 0$ .

\vfill

\newpage

\textbf{1-9}\quad 求一维氢原子定态能量和波函数, $\displaystyle V(x) = -\dfrac{e^2}{|x|}$ .

\vfill

\newpage

\textbf{1-10}\quad 粒子在一维势场 $\displaystyle V(x) = V_{0}\left(\dfrac{a}{x} -\dfrac{x}{a}\right)^{2}$ 中运动, 其中 $\displaystyle V_{0}, a$ 是正的常数. 求定态能量和波函数.

\vfill

\newpage

\textbf{1-11}\quad 设粒子的波函数 $\displaystyle \psi(x) = A\left(\dfrac{x}{a}\right)^n\mathrm{e}^{-x / a}$ 是一维势 $\displaystyle V(x)$ 中的粒子能量本征态, 其中 $\displaystyle A, a$ 和 $\displaystyle n$ 为常数. 当 $\displaystyle x \to \infty$ 时, $\displaystyle V(x) \to 0$ . 试求势能 $\displaystyle V(x)$ 和粒子能量 $\displaystyle E$ . 粒子质量为 $\displaystyle \mu$ .

\vfill

\newpage

\textbf{1-12}\quad 已知一维定态波函数为 $\displaystyle \psi(x) = \begin{cases} a^2 - x^2, & |x| < a \\ 0, & |x| > a \end{cases}$ , 且有 $\displaystyle \langle \psi | V | \psi \rangle = 0$ . 试从一维定态方程出发, 求出势位函数 $\displaystyle V(x)$ 与定态能量 $\displaystyle E$ .

\vfill

\newpage

\textbf{1-13}\quad 质量为 $\displaystyle \mu$ 的粒子在势场 $\displaystyle V(x) = \begin{cases} \infty, & x < 0 \\ 0, & 0 \leqslant x \leqslant a \\ V_0, & x > a \end{cases}$ 中运动 $\displaystyle (V_0 > 0)$ . 已知该粒子在此势场中存在一能量为 $\displaystyle E = V_0 / 2$ 的态. 试确定此势阱的宽度 $\displaystyle a$ .

\vfill

\newpage

\textbf{1-14}\quad 质量为 $\displaystyle \mu$ 的粒子在势场 $\displaystyle V(x)$ 中作一维束缚运动, 两个能量本征函数为 $\displaystyle \psi_{1}(x) = A\mathrm{e}^{-\beta x^{2} / 2}, \psi_{2}(x) = B(x^{2} + bx + c)\mathrm{e}^{-\beta x^{2} / 2}$ , $\displaystyle A, B, b, c$ 均为实常数. 试确定参数 $\displaystyle b, c$ 的取值, 并求这两个态的能量之差 $\displaystyle E_{2} - E_{1}$ .

\vfill

\newpage

\textbf{1-15}\quad 质量为 $\displaystyle \mu$ 的粒子在一圆圈(周长为 $\displaystyle L$ )上运动, 并受到 $\displaystyle \delta$ 函数势 $\displaystyle V(x) = a\delta (x - (L / 2))$ 的作用. 求粒子能级和相应的归一化波函数.

\vfill

\newpage

\textbf{1-16}\quad 粒子在势场 $\displaystyle V(x)$ 中作一维运动， $\displaystyle \hat{H}_{0}=\dfrac{\hat{p}^{2}}{2\mu}+V(x)$ ，定态能量为 $\displaystyle E_{n}^{(0)}, n=1,2,\cdots$ 。求哈密顿量 $\displaystyle \hat{H}=\hat{H}_{0}+\dfrac{\lambda}{2}\hat{p}$ 的本征值， $\displaystyle \lambda$ 为参数。

\vfill

\newpage

\textbf{1-17}\quad 电荷为 $\displaystyle q$ 质量为 $\displaystyle \mu$ 的点粒子在一维均匀电场 $\displaystyle E$ 中运动, 位势为 $\displaystyle V(x) = -qEx$ . 在 $\displaystyle t = 0$ 时该粒子的坐标与动量平均值分别为 $\displaystyle \langle x \rangle = x_0$ 与 $\displaystyle \left\langle \hat{p}_x \right\rangle = 0$ . (1) 计算 $\displaystyle t$ 时刻的动量平均值 $\displaystyle \left\langle \hat{p}_x \right\rangle$ ; (2) 计算 $\displaystyle t$ 时刻的坐标平均值 $\displaystyle \langle x \rangle$ ; (3) 把计算的结果同经典物理的结果比较.

\vfill

\newpage

\textbf{1-18}\quad 质量为 $\displaystyle \mu$ 的粒子被约束在半径为 $\displaystyle r$ 的圆周上运动。(1)设立路障, 进一步限制粒子在 $\displaystyle 0 < \varphi < \varphi_0$ 的一段圆弧上运动, $\displaystyle V(\varphi) = \begin{cases} 0, & 0 \leqslant \varphi < \varphi_0 \\ \infty, & \varphi_0 \leqslant \varphi < 2\pi \end{cases}$ , 求解粒子的本征能量和本征函数; (2)设粒子处于情况(1)的基态, 求突然撤去路障后, 粒子仍处于最低能量态的概率.

\vfill

\newpage

\textbf{1-19}\quad 质量为 $\displaystyle \mu$ 的粒子处于一维谐振子势场 $\displaystyle V_{1} = \dfrac{1}{2} kx^{2}$ 的基态, 某时刻弹性系数 $\displaystyle k$ 突然变为 $\displaystyle 2k$ , 即势场变为 $\displaystyle V_{2} = kx^{2}$ . 求此时粒子处于新势场 $\displaystyle V_{2}$ 的基态的概率, 并求此后粒子能量的平均值.

\vfill

\newpage

\textbf{1-20}\quad 同上题. 当 $\displaystyle k$ 变成 $\displaystyle 2k$ 后, 经过多少时间 $\displaystyle T$ 再将 $\displaystyle 2k$ 变为 $\displaystyle k$ , 粒子 $\displaystyle 100\%$ 回到原来的基态?

\vfill

\newpage

\textbf{1-21}\quad 质量为 $\displaystyle \mu$ 的粒子处于一维刚性盒 $\displaystyle (0 \sim a)$ 的基态. 盒子的 $\displaystyle x = a$ 壁突然运动至 $\displaystyle x = 2a$ 处, 试计算盒子膨胀后粒子仍处于基态的概率. 如果盒子的两壁对称地向两边移动, 盒子的宽度由 $\displaystyle a$ 变为 $\displaystyle 2a$ , 结果又如何?

\vfill

\newpage

\textbf{1-22}\quad 一个粒子处于宽度为 $\displaystyle a$ 的无限深方势阱中的基态. 若(1)阱的两壁同时缓慢地由宽度 $\displaystyle a$ 缩小为 $\displaystyle a / 2$ ; (2)阱的两壁同时突然地由宽度 $\displaystyle a$ 缩小为 $\displaystyle a / 2$ , 求粒子留在基态的概率.

\vfill

\newpage

\textbf{1-23}\quad 一个粒子处于宽度为 a 的无限深方势阱的基态，能量为 $\displaystyle E_{1}=38eV$ .

(1) 计算第一激发态能量；

(2) 计算基态粒子对阱壁的平均作用力.

\vfill

\newpage

\textbf{1-24}\quad 质量为 $\displaystyle \mu$ 的粒子处于 $\displaystyle 0 \leqslant x \leqslant a$ 的无限深方势阱中. $\displaystyle t = 0$ 时, 归一化波函数为 $\displaystyle \psi(x,0) = \sqrt{\dfrac{8}{5a}}\left(1 + \cos \dfrac{\pi x}{a}\right)\sin \dfrac{\pi x}{a}$ . 求

(1)在后来某一时刻 $\displaystyle t_0$ 的波函数;

(2)在 $\displaystyle t = 0$ 与 $\displaystyle t = t_0$ 时体系的能量;

(3)在 $\displaystyle t_0$ 时粒子处于 $\displaystyle 0 \leqslant x \leqslant a / 2$ 内的概率.

\vfill

\newpage

\textbf{1-25}\quad 质量为 $\displaystyle \mu$ 的粒子在无限深方势阱 $\displaystyle V(x) = \begin{cases} \infty, & x < 0, x > a \\ 0, & 0 < x < a \end{cases}$ 中运动. $\displaystyle t = 0$ 时粒子处于状态 $\displaystyle \psi(x,0) = A\sin \dfrac{\pi x}{2a}\cos^3\dfrac{\pi x}{2a}$ , 其中 $\displaystyle A$ 为常数. 求出 $\displaystyle t$ 时刻粒子,

(1) 处于基态的概率;

(2) 能量平均值;

(3) 动量平均值;

(4) 动量均方差根值(不确定度).

\vfill

\newpage

\textbf{1-26}\quad 质量为 $\displaystyle m$ 的粒子位于一维无限深方势阱 $\displaystyle [-a / 2, a / 2]$ 中, 势阱宽度为 $\displaystyle a$ . $\displaystyle t = 0$ 时体系处于无限深方势阱中能量最低的两个态的线性叠加态, 各自的概率为 $\displaystyle 50\%$ . 计算 $\displaystyle t > 0$ 时粒子的概率密度和动量平均值.

\vfill

\newpage

\textbf{1-27}\quad 已知 $\displaystyle t = 0$ 时一维运动粒子在态 $\displaystyle \psi(x)$ 中坐标 $\displaystyle x$ 和动量 $\displaystyle \hat{p}$ 的平均值分别为 $\displaystyle x_0$ 与 $\displaystyle p_0$ , 求 $\displaystyle t = 0$ 时在态 $\displaystyle \varphi(x) = e^{-ip_0x / \hbar} \psi(x + x_0)$ 中 $\displaystyle x$ 与 $\displaystyle \hat{p}$ 的平均值.

\vfill

\newpage

\textbf{1-28}\quad 粒子处于宽度为 $\displaystyle a$ 的一维无限深方势阱 $\displaystyle V(x) = \begin{cases} 0, & 0 < x < a \\ \infty, & x < 0, x > a \end{cases}$ 中的定态 $\displaystyle \psi_n(x)$ , 求粒子的动量分布概率 $\displaystyle \left|\varphi(p)\right|^2$ .

\vfill

\newpage

\textbf{1-29}\quad 粒子处于宽度为 $\displaystyle a$ 的一维无限深方势阱的基态, $\displaystyle t = 0$ 时阱的两壁突然崩溃. 求 $\displaystyle t > 0$ 时粒子处于动量取值在 $\displaystyle p \sim p + \mathrm{d}p$ 内的概率, 以及粒子波函数的表示式(不必算出结果).

\vfill

\newpage

\textbf{1-30}\quad 质量为 $\displaystyle \mu$ 的粒子在一维势场 $\displaystyle V(x) = \begin{cases} 0, & |x| < a \\ \infty, & |x| > a \end{cases}$ 中运动。

(1)求粒子定态能量 $\displaystyle E_{n}$ 与归一化定态波函数 $\displaystyle \psi_{n}(x)$ .

(2)求粒子在定态 $\displaystyle \psi_{n}(x)$ 上的平均值 $\displaystyle \overline{x}$ .

(3)设 $\displaystyle t = 0$ 时粒子波函数为 $\displaystyle \psi(x,0) = \begin{cases} A(a^2 - x^2), & |x| < a \\ 0, & |x| > a \end{cases}$ , 其中 $\displaystyle A$ 为归一化常数, 求

$\displaystyle (a)$ 在 $\displaystyle \psi(x,0)$ 态上粒子能量取值 $\displaystyle E_{n}$ 的概率;

(b)粒子的平均能量 $\displaystyle \overline{E}$ ;

(c)任意 $\displaystyle t > 0$ 时粒子波函数 $\displaystyle \psi(x,t)$ 表示式.

\vfill

\newpage

\textbf{1-31}\quad 设一维运动粒子的坐标和动量分别为 $\displaystyle q$ 和 $\displaystyle \hat{p}$ .

(1) 计算 $\displaystyle \hat{p}$ 和 $\displaystyle \mathrm{e}^{\mathrm{i}cq}$ 的对易关系, 其中 $\displaystyle c$ 为常数;

(2) 若 $\displaystyle p_0$ 是 $\displaystyle \hat{p}$ 的本征值, 相应的本征函数是 $\displaystyle \psi_0(q)$ , 证明 $\displaystyle p_0 + \hbar c$ 也是 $\displaystyle \hat{p}$ 的本征值, 给出相应的本征函数.

\vfill

\newpage

\textbf{1-32}\quad 写出二维谐振子势 $\displaystyle V = \dfrac{1}{2}\mu \omega_x^2 x^2 +\dfrac{1}{2}\mu \omega_y^2 y^2$ 中粒子能级：

(1)设 $\displaystyle \omega_{x} = \omega_{y}$ 求能级简并度；

(2)设 $\displaystyle \omega_{x} / \omega_{y} = \dfrac{1}{2}$ ，求能级简并度.

\vfill

\newpage

\textbf{1-33}\quad 粒子在二维势场 $\displaystyle V(x, y) = \dfrac{1}{2} \mu \omega^2 (x^2 + y^2 + 2\lambda xy)$ 中运动, 其中 $\displaystyle \left|\lambda\right| < 1$ , $\displaystyle \mu$ 为粒子质量. 求能量本征值和本征函数.

\vfill

\newpage

\textbf{1-34}\quad 两个质量都是 $\displaystyle \mu$ 的一维耦合谐振子体系的哈密顿量为
$$
\hat {H} = \dfrac {1}{2 \mu} (\hat {p} _ {1} ^ {2} + \hat {p} _ {2} ^ {2}) + \dfrac {1}{2} \mu \omega^ {2} [ (x _ {1} - a) ^ {2} + (x _ {2} + a) ^ {2} + \lambda (x _ {1} - x _ {2}) ^ {2} ]
$$
其中 $\displaystyle \lambda$ 与 a 为参数, $\displaystyle \lambda > -\dfrac{1}{2}, -\infty < x_{1}, x_{2} < \infty$ . 求体系能量本征值.

\vfill

\newpage

\textbf{1-35}\quad 质量为 $\displaystyle \mu$ 的粒子在势场
$$
V (x, y, z) = A \left(x ^ {2} + y ^ {2} + 2 \lambda x y\right) + B \left(z ^ {2} + 2 c z\right)
$$
中运动, $\displaystyle A, B > 0, |\lambda| < 1$ , $\displaystyle c$ 取任意实数, 求能量本征值. 如考虑另一个新势 $\displaystyle V'$ , 它同原势的关系为 $\displaystyle V' = \begin{cases} V, & z > -c, xy \text{任意} \\ \infty, & z < -c, xy \text{任意} \end{cases}$ , 求基态能量.

\vfill

\newpage

\textbf{1-36}\quad 设一维粒子由 $\displaystyle x = -\infty$ 处以平面波 $\displaystyle \psi_{in} = e^{ikx}$ 入射，在原点处受到势能 $\displaystyle V(x) = V_0\delta (x)$ 的作用.

(1)写出波函数的一般表达式；

(2)确定粒子波函数在原点处满足的边界条件；

(3)求出该粒子的透射系数和反射系数；

(4)分别指出 $\displaystyle V_{0} > 0$ 与 $\displaystyle V_{0} <   0$ 时的量子力学效应.

\vfill

\newpage

\textbf{1-37}\quad 在以下两种情况中计算入射粒子在一维阶跃势上的反射率 $\displaystyle R$ 与透射率 $\displaystyle T, V(x) = \begin{cases} 0, & x < 0 \\ V_0, & x > 0 \end{cases}$ :

(1) $\displaystyle E > V_0$ ;

(2) $\displaystyle E < V_0$ .

\vfill

\newpage

\textbf{1-38}\quad 电子经 1000V 电压加速后由 $\displaystyle x = -\infty$ 射向阶跃势 $\displaystyle V(x) = \begin{cases} 0, & x < 0 \\ V_0, & x > 0 \end{cases}$ , 其中 $\displaystyle V_0 = 750\mathrm{eV}$ . 现有 1800 个电子入射, 在 $\displaystyle x = \infty$ 处能观测到多少个电子?

\vfill

\newpage

\textbf{1-39}\quad 粒子被一维矩形势垒 $\displaystyle V(x) = \begin{cases} 0, & x < 0, x > a \\ V_0, & 0 < x < a \end{cases}$ 散射。

(1)当粒子的能量 $\displaystyle E > V_0$ 时，求反射率 $\displaystyle R$ 与透射率 $\displaystyle T$ ；

(2)当粒子的能量 $\displaystyle E = V_0$ 时，有一半粒子被反射回去，求粒子的质量所满足的方程。

\vfill

\newpage

\textbf{1-40}\quad 把传导电子限制在金属内部的是一种平均势, 对于下面的一维势模型: $\displaystyle V(x) = \begin{cases} -V_0, & x < 0 \\ 0, & x > 0 \end{cases}$ , 试就

(1) $\displaystyle E > 0$ ;

(2) $\displaystyle -V_0 < E < 0$ 两种情况计算接近金属表面的传导电子的反射率与透射率.

\vfill

\newpage

\textbf{1-41}\quad 质量为 $\displaystyle m$ 的电子以动能 $\displaystyle E > V_{0}$ 由左向右入射到高度为 $\displaystyle V_{0}$ 的台阶势上, 在台阶势的跃起处还存在 $\displaystyle \delta$ 势: $\displaystyle \gamma \delta(x)(\gamma > 0)$ , 即考虑电子在势 $\displaystyle V(x) = V_{0} \theta(x) + \gamma \delta(x), \theta(x) = \begin{cases} 0, & x < 0 \\ 1, & x > 0 \end{cases}$ , 上的散射, $\displaystyle \theta(x)$ 为单位阶跃函数.

(1) 列出定态薛定谔方程及波函数导数 $\displaystyle \psi'$ 在 $\displaystyle x = 0$ 两侧的跃变条件;

(2) 求电子在 $\displaystyle x = 0$ 处的透射系数 $\displaystyle T = \left| \dfrac{j_{out}}{j_{in}} \right|$ 和反射系数 $\displaystyle R = \left| \dfrac{j_{ref}}{j_{in}} \right|$ .

\vfill

\newpage

\textbf{1-42}\quad 求一维常虚势场 $\displaystyle -\mathrm{i}V(V \ll E)$ 中运动的粒子波函数, 计算概率流密度并证明虚势代表粒子被吸收, 求吸收系数 $\displaystyle M$ , 用 $\displaystyle V$ 表示.

\vfill

\newpage

\textbf{1-43}\quad 质量为 $\displaystyle \mu$ 的粒子在势场 $\displaystyle V(x) = \begin{cases} \infty, & x < 0 \\ -V_0 a \delta (x - a), & x > 0 \end{cases}$ 中运动, 其中 $\displaystyle V_{0}$ 和 $\displaystyle a$ 都是正实数. 求

(1)束缚态能量满足的方程;

(2)存在束缚态的条件.

\vfill

\newpage

\textbf{1-44}\quad 质量为 $\displaystyle \mu$ 的粒子在势场 $\displaystyle V(x) = \begin{cases} 0, & 0 < x < a \\ \infty, & x < 0, x > a \end{cases}$ 中运动. $\displaystyle t = 0$ 时粒子的波函数为 $\displaystyle \psi(x,0) = A\left(1 + 2b\cos \dfrac{\pi x}{a}\right)\sin \dfrac{\pi x}{a}$ , 其中 $\displaystyle A, b$ 为常数. 求任意 $\displaystyle t$ 时粒子的波函数 $\displaystyle \psi(x,t)$ , 平均能量 $\displaystyle \overline{E(t)}$ 和平均动量 $\displaystyle \overline{p(t)}$ .

\vfill

\newpage

\textbf{1-45}\quad 对于一维谐振子哈密顿量 $\displaystyle \hat{H} = \dfrac{\hat{p}^2}{2\mu} + \dfrac{1}{2}\mu \omega^2 x^2$ ，求海森伯绘景中的坐标 $\displaystyle \hat{x}(t)$ 与动量 $\displaystyle \hat{p}(t)$ .

\vfill

\newpage

\textbf{1-46}\quad 若在薛定谔绘景中, 体系的哈密顿量 $\displaystyle \hat{H} = \omega \hat{L}_z$ , 试给出在海森伯绘景中的 $\displaystyle \hat{L}_x(t)$ 与 $\displaystyle \hat{L}_y(t)$ .

\vfill

\newpage

\textbf{1-47}\quad 设一维体系能量算符 $\displaystyle \hat{H} = \dfrac{1}{2\mu}\left(\hat{p}^2 - \dfrac{\hbar^2}{x^2}\right)$ .

(1) 利用维里定理证明该体系不存在束缚态；

(2) 用海森伯运动方程证明算符 $\displaystyle \hat{Q}(t) = \dfrac{1}{4}(x\hat{p} + \hat{p}x) - \hat{H}t$ 在任意态上的平均值为常数.

\vfill

\newpage

\textbf{1-48}\quad 什么是量子化？如何实现量子化？

\vfill

\newpage

